\documentclass[11pt]{article}

\usepackage[T1]{fontenc}
\usepackage{lmodern}
\usepackage[a4paper,margin=26mm]{geometry}
\usepackage{amsmath,amssymb,amsthm,mathtools}
\usepackage{microtype}
\usepackage{enumitem}
\usepackage{booktabs,array,tabularx,longtable}
\usepackage{xcolor}
\usepackage[colorlinks=true,linkcolor=blue!48!black,
  citecolor=blue!48!black,urlcolor=blue!55!black]{hyperref}
\setlength{\emergencystretch}{2em}

\newtheorem{theorem}{Theorem}[section]
\newtheorem{proposition}[theorem]{Proposition}
\newtheorem{corollary}[theorem]{Corollary}
\newtheorem{lemma}[theorem]{Lemma}
\theoremstyle{definition}
\newtheorem{definition}[theorem]{Definition}
\theoremstyle{remark}
\newtheorem{remark}[theorem]{Remark}

\newcommand{\C}{\mathbb C}
\newcommand{\R}{\mathbb R}
\newcommand{\N}{\mathbb N}
\newcommand{\cA}{\mathcal A}
\newcommand{\cD}{\mathcal D}
\newcommand{\cE}{\mathcal E}
\newcommand{\cF}{\mathcal F}
\newcommand{\cG}{\mathcal G}
\newcommand{\cH}{\mathcal H}
\newcommand{\cI}{\mathcal I}
\newcommand{\cK}{\mathcal K}
\newcommand{\cQ}{\mathcal Q}
\newcommand{\cR}{\mathcal R}
\newcommand{\cS}{\mathcal S}
\newcommand{\cU}{\mathcal U}
\newcommand{\cW}{\mathcal W}
\newcommand{\Rea}{\operatorname{Re}}
\newcommand{\im}{\operatorname{im}}
\newcommand{\kerop}{\operatorname{ker}}
\newcommand{\ran}{\operatorname{ran}}
\newcommand{\Tr}{\operatorname{Tr}}
\newcommand{\Div}{\operatorname{Div}}
\newcommand{\Dom}{\operatorname{Dom}}
\newcommand{\cl}{\operatorname{cl}}

\title{\textbf{Relative admission spectra at the Riemann divisor}\\
\large A conservation law for arithmetic sources and singular proof transports}
\author{Michael Haufschild}
\date{}

\begin{document}

\maketitle

\begin{abstract}
Arithmetic approaches to the Riemann divisor now possess several genuine
sources and several nonisomorphic transports.  The unresolved problem is
therefore not the absence of a source, but deciding which part of a source
survives passage to a zero-excluding boundary condition.  We introduce a
route invariant which answers this question.

For a differential \(d:A_0\to A_1\), an admitted primitive sector
\(S\subset A_0\), a primitive selector \(h\), and a parameterized source
family \(x_p\), define
\[
 \mathfrak o_S(p)=[h(x_p)]\in A_0/(S+\ker d),\qquad
 \Sigma_S(x)=\{p:\mathfrak o_S(p)=0\}.
\]
We prove the relative admission-spectrum conservation law: the class is
independent of the primitive selector, its vanishing is equivalent to the
existence of an admitted primitive, and it is natural under every chain
transport preserving admitted sectors.  An injective induced quotient
transport preserves the spectrum exactly.  If a target admits a parameter
which the source excludes, the transport has killed a specific nonzero
relative class.  Thus spectral mismatch forces information loss or a
singular change of domain.

The law applies without identifying the underlying route categories.  In
the all-prime Calder\'on--Witt complex it recovers Green-incidence
transgression.  In two independently constructed analytic sources it gives
a sharp fracture: the Hedenmalm Borel reflection family is square
integrable throughout the centered open half-strip, whereas the polarized
Euler--Hardy reflection family is square integrable only at zero modular
displacement.  Consequently no linear map preserving even the one-way
square-integrable domain can intertwine the two families at a nonzero strip
parameter.  Any bridge must be domain-singular; bounded-operator completion
cannot supply it.

We also consolidate the constructed source prefix.  Target-free
M\"obius--theta translates have Fourier--Laplace transform
\(\zeta(ms)^{-1}\Xi\) on their absolute-convergence domain, while actual
Hellinger--Beta residues lift to nonzero Calder\'on--Witt incidences whose
relative Gamma boundary class vanishes exactly at zero displacement.
These results leave one precise missing producer: a source-derived,
off-diagonal cofinal boundary transport which carries the physical source
into the sharp Euler--Hodge admission sector.  We do not prove that landing
or the Riemann hypothesis; we prove why no regular identification of the
existing source carriers can replace it.
\end{abstract}

\medskip
\noindent\textbf{Keywords.}
Riemann hypothesis; admission spectrum; source transport; Green operator;
Calder\'on--Witt incidence; Hedenmalm space; Gamma/co-Poisson transform;
Hodge polarization.

\medskip
\noindent\textbf{2020 Mathematics Subject Classification.}
Primary 11M26, 11M06; Secondary 46N50, 47A06, 58A14.

\section{Introduction}

The research frontier around a source-based proof of the Riemann
hypothesis is easily misdescribed as a search for one missing source or
one missing transport.  Neither description is accurate.  Arithmetic,
analytic, cohomological, probabilistic, spectral, and recurrence sources
have been constructed.  They do not have the same finite-energy domain,
the same primitive sector, or the same notion of faithful continuation.
Their coexistence creates several live routes, not several names for one
operator problem.

The apparent multiplicity of proof routes creates a second danger.  An
RH-equivalent positivity, compactness, self-adjointness, recurrence, or
growth condition is not by itself an independent route.  It is a
\emph{consumer}: a representation of zero displacement after some source
and transport have already been chosen.  It becomes a proof route only
when an unconditional producer theorem forces that consumer condition.

The useful route-independent decomposition is
\[
 \boxed{
 \text{source}\ \longrightarrow\
 \text{faithful transport}\ \longrightarrow\
 \text{oriented admission}\ \longrightarrow\
 \text{consumer}.}
\tag{1.1}\label{eq:four-stage}
\]
Different constructions stop at different arrows.  The new issue is to
compare them without pretending that their source spaces coincide.  We do
so by forgetting the carrier but retaining the set of parameters for which
its source primitive reaches the route's admitted sector.  This
\emph{relative admission spectrum} is conserved by every source-faithful
chain transport.  A mismatch between two spectra is therefore an exact
obstruction to regular transport, not a failure to guess the right kernel.

The common RH-bearing parameter is the real modular displacement
\[
                    \delta_\rho
                    =2\left(\Rea\rho-\frac12\right).
\tag{1.2}\label{eq:delta}
\]
It is retained before scalarization as a source-tail entropy, a
Cayley-horizon escape rate, a Frobenius--Hodge commutator, a reflected
valuation type, a Gamma sewing exponent, and the horizontal part of the
completed logarithmic-current residue.  The problem is not to rediscover
\(\delta_\rho\).  It is to force a source relation carrying
\(\delta_\rho\) into a boundary class which is admitted only when
\(\delta_\rho=0\).

Earlier source laws established three hard boundaries.
\begin{enumerate}[label=\textup{(\roman*)},leftmargin=2.2em]
\item Complete prime-normal retention in the product topology makes a
      closed relation maximally vertical and therefore output-free.
\item Reflected valuation sectors at nonzero displacement have
      inequivalent asymptotic type; exact sewing grows beyond every finite
      Rees order.
\item Local entire carriers, ordinary closable positive source metrics,
      bounded Ward windows, and finite differential-algebraic closures
      cannot retain the required oriented charge.
\end{enumerate}
These results say what the terminal arrow in \eqref{eq:four-stage} cannot
be.  They do not say that transport is absent.

The decisive organizing input is relative admission-spectrum
conservation.  The algebraic source has one absolute orientation line;
every other invariant is the relative escape of a Green primitive from a
source-selected admitted sector.  Primitive choice cannot change the
class.  Faithful transport cannot change its vanishing set.  This explains
why regular scalar consumers repeatedly lose the sign and why singular,
relation-valued, or source-domain constructions recur across otherwise
different programs.

Two recent constructions sharpen this statement.  First, physical
M\"obius--theta translates now land in the common Xi divisor before a zero
is selected.  Second, Hellinger--Beta residues now land in actual
cross-prime Calder\'on--Witt incidence and detect every nonzero modular
displacement in the relative Gamma boundary.  These close source and
finite-witness gaps, but their admission spectra do not match under a
regular Hilbert transport.  The surviving bridge must change domain and
polarization cofinally.

The paper has one main law and three independent instantiations.
Sections~\ref{sec:sources} and~\ref{sec:transport} inventory constructed
producers and transports.  Section~\ref{sec:admission} proves the relative
admission-spectrum law, the all-prime incidence specialization, and the
Hedenmalm--Euler fracture.
Section~\ref{sec:routes} identifies the precise obligation of each
remaining route.  Sections~\ref{sec:shadows} and~\ref{sec:proof} explain
why the no-go theorems and positive consumers are shadows of the
same law and state the exact theorem which would complete the RH proof.

\section{The common oriented invariant}
\label{sec:charge}

Let \(\Xi\) denote the completed Riemann function and let \(\rho\) be a
nontrivial divisor point of multiplicity \(m_\rho\).  Its centered outer
charge is
\[
                    \chi_\rho
                    =-m_\rho\left(\rho-\frac12\right).
\tag{2.1}\label{eq:outer-charge}
\]
The horizontal charge is
\[
             \Rea\chi_\rho
              =-\frac{m_\rho}{2}\,\delta_\rho .
\tag{2.2}\label{eq:horizontal-charge}
\]
This is not a coordinate chosen after the divisor is known.  It has
several exact source realizations.

\subsection{Prime-tail entropy and source-horizon rate}

Let \(B(s)\) be the canonical prime-two shell-transfer character and
define
\[
 \mathsf E(s)=
 \frac{\log\lvert B(s)\rvert-
       \log\lvert B(1-\overline s)\rvert}{\log2}.
\tag{2.3}\label{eq:entropy}
\]
Then
\[
 \mathsf E(s)=-2\left(\Rea s-\frac12\right)
             =-\delta_s.
\tag{2.4}\label{eq:entropy-delta}
\]
The same number is the negative Cayley source-horizon escape rate.
After multiplication by \(m_\rho/2\), it gives
\(\Rea\chi_\rho\).  Two prime clocks determine the full complex charge
\(\chi_\rho\), not merely its real part.

Let \(C_{m,\rho}\) be the source Cayley residue mode with pole multiplicity
\(m\), and let \(x=x_\rho(\tau)\) be its canonical observation point.
Every order-\(r\) jet then satisfies the exact Poisson conservation
identity
\[
 \frac{x^r}{r!}
 \left\lvert C_{m,\rho}^{(r)}(x)\right\rvert^2
 =
 |m|^2 e^{-\mathsf E(\rho)\tau}
 \frac{e^{-\tau}\tau^r}{r!},
\tag{2.5}\label{eq:poisson-conservation}
\]
so the critical line is the neutral horizon
\(\mathsf E(\rho)=0\).

\subsection{The same charge in boundary coordinates}

The nearby-cycle duality--adjunction fracture is
\[
 \mathfrak f(\rho)=2\left(\Rea\rho-\frac12\right)
                  =\delta_\rho.
\tag{2.6}\label{eq:fracture}
\]
Its vanishing is equivalent to each of the following exact conditions:
\begin{enumerate}[label=\textup{(\alph*)},leftmargin=2.2em]
\item weight-one Frobenius--Hodge commutation;
\item neutrality of the source modular character;
\item equality of reflected prime-seam charges;
\item zero reflected valuation-entropy difference;
\item unit Gamma/co-Poisson Hodge sewing;
\item two-face reflection equipartition.
\end{enumerate}
These are coordinate identities for one source invariant.  Their logical
equivalence does not prove its vanishing.

\section{Constructed sources}
\label{sec:sources}

The constructions supply several nonisomorphic sources.  Their coexistence is
important: a source failure in one category does not imply absence of
arithmetic data.

\subsection{Mass-one prime-normal resolution}

For \(K\ge12\), define
\[
 \Phi^K_0(s)=\zeta(Ks)^{-1},\qquad
 \Phi^K_{j+1}(s)=
   \zeta((K+j+1)s)^{-1}-\zeta((K+j)s)^{-1}.
\tag{3.1}\label{eq:positive-resolution}
\]
The vector \(E_K(s)=(\Phi^K_j(s))_{j\ge0}\) is a mass-one,
coefficient-disjoint, divisor-faithful source on
\(\Rea s>\tfrac12\).  It is injective for every finite \(K\), while
\(E_K(s)\to e_0\) in \(\ell^1\).  Exact shell peeling recovers the
countable normal reader
\[
                        \bigl(p^{1/2-s}\bigr)_p .
\tag{3.2}\label{eq:normal-reader}
\]
At \(s_j=\tfrac12+j\), every finite leading prime block is a nonsingular
Vandermonde matrix.

\subsection{Physical M\"obius--theta source and Xi landing}

Let \(\Phi\) be the full theta-cusp source, whose Fourier--Laplace
transform is the centered completed function \(\Xi_C\) on
\(|\Im w|<1\), and put
\[
                         s(w)=\frac12+iw.
\tag{3.3}\label{eq:centered-parameter}
\]
For an integer \(m\ge3\), the power-supported M\"obius source rung is the
target-free translate sum
\[
 \Psi_m(x)=\sum_{n\ge1}
   a_m(n)n^{-1/2}\Phi(x-\log n),
 \qquad a_m=\operatorname{mPart}_m(\mu).
\tag{3.4}\label{eq:physical-mobius-theta}
\]
Thus \(a_m(k^m)=\mu(k)\) and \(a_m(n)=0\) off the set of perfect
\(m\)-th powers.
The corresponding series converges in the physical \(L^2\) source.  On
the joint absolute-convergence domain
\[
            |\Im w|<1,\qquad \Rea s(w)>\frac1m,
\]
termwise translation and Bochner integration give the exact landing
\[
 \int_{\R}e^{-iwx}\Psi_m(x)\,dx
   =\zeta\!\left(ms(w)\right)^{-1}\Xi_C(w).
\tag{3.5}\label{eq:physical-xi-landing}
\]
The source in \eqref{eq:physical-mobius-theta} contains no spectral point,
zero, or value of \(\Xi_C\).  The common divisor appears only after the
source transform is applied.

For \(K\ge3\), define a complete source channel before transform by
\[
 C_K(0)=\Psi_K,\qquad
 C_K(j+1)=\Psi_{K+j+1}-\Psi_{K+j}.
\tag{3.6}\label{eq:complete-physical-channel}
\]
Every coordinate has an actual Fourier--Laplace integral and
\[
 \mathcal F_L C_K(j)(w)
   =M_{K,j}(s(w))\Xi_C(w),
\tag{3.7}\label{eq:complete-physical-landing}
\]
where \(M_{K,j}\) is the head/difference reciprocal-zeta multiplier.
Thus Route~R now has a physical source-to-Xi producer on its honest
complex domain.  Its open theorem begins after this landing: the source
must be transported into a sharp polarized admission domain without
losing its prime-normal labels.

\subsection{Valuation probability and polarized boundary}

Diagonal Euler activation is
\[
                          \frac{\zeta(s)}{\zeta(2s)}.
\tag{3.8}\label{eq:activation}
\]
At \(q=p^{-\sigma}\), its local valuation law is
\[
 \nu_q(0)=\frac1{1+q},\qquad
 \nu_q(k)=\frac{(1-q)q^k}{1+q}\quad(k\ge1).
\tag{3.9}\label{eq:valuation-law}
\]
It disintegrates into Bernoulli squarefree occupancy followed by a
conditional geometric valuation fibre.  On the continuous source
boundary,
\[
 h(t)=e^{-t},\qquad \tau(t)=1-e^{-t}.
\tag{3.10}\label{eq:head-tail}
\]
The forced odd Hodge scale
\[
 q_H(t)^2=\frac{\tau(t)}{h(t)}
\tag{3.11}\label{eq:hodge-scale}
\]
is reciprocal Fisher whitening.  Both cross-channels equal
\(\sqrt{h(t)\tau(t)}\), the Hellinger and KMS/Wightman half-density.
Its common-test transform is the exact Beta/Gamma column
\[
 -\int_0^\infty e^{-at}\sqrt{h(t)\tau(t)}\,dt
 =-\frac{\Gamma(a+\tfrac12)\Gamma(\tfrac32)}
          {\Gamma(a+2)}.
\tag{3.12}\label{eq:beta-gamma}
\]

The source also contains a transverse score.  If
\(S_0(t)=2(1-t)\), the residual
\[
 X_\perp(t)=\bigl(\tau(t)S_0(t),-h(t)S_0(t)\bigr)
\tag{3.13}\label{eq:transverse}
\]
is orthogonal to every diagonal common Laplace test and has squared
polarized norm \(3\).  The canonical normalization
\((2/\sqrt3)X_\perp\) has squared norm \(4\).  Thus the four-unit
orientation is source-generated; its identification with an actual Xi
Ward column remains a transport statement.

\subsection{Jet, crease, theta, and relation sources}

Four further sources play distinct roles.
\begin{enumerate}[leftmargin=2.2em]
\item The multiplicity-exact Xi jet source consists of finite divisor jets
      and their \(\ell^2\) completion.  Its compact-support Green map is
      continuous and dense.
\item The first CCM clock has derivative creases
      \[
      D_+T(\log k)-D_-T(\log k)
       =-\frac{2\pi}{\log k}\,k^{-1/2}\Lambda(k),
      \tag{3.14}\label{eq:ccm-crease}
      \]
      and therefore recovers every von Mangoldt atom.
\item The theta completed-density source carries, at every Xi zero, an
      exact reflected zero relation.  A Hurwitz functional-equation pair
      carries the corresponding differential circuit exactly on the
      critical line.
\item Universal facet sources are differentially independent, have a
      countable point-finite differential matroid, and carry summable
      residue-orientation weights.  Projective recurrence
      sources encode exact balanced spectral relations as the kernel of
      one energy detuning map.
\end{enumerate}
These sources are not interchangeable.  They become comparable only
after specifying what relation their transport must conserve.

\section{Constructed transports and hard exclusions}
\label{sec:transport}

\subsection{The product-topology and finite-Rees boundaries}

The positive resolution \eqref{eq:positive-resolution} collapses
horizontally while retaining every finite normal block.  If a closed
linear relation in the product target retains those blocks, its
multivalued part is dense and hence equal to the whole target:
\[
 \operatorname{mul}\cR=\C^\N,\qquad
 \cR=\Dom(\cR)\times\C^\N.
\tag{4.1}\label{eq:maximal-vertical}
\]
It transports no output information.  Therefore the relevant completion
cannot use the ordinary countable product topology.

Independently, at any nonzero displacement one reflected valuation sector
is finite or half-normal while the other is flat beyond every finite
Rees order.  If \(R_N\) sews them exactly, then for every \(K\),
\[
 a_N^{-K}=o(\|R_N\|).
\tag{4.2}\label{eq:beyond-rees}
\]
Thus an informative source completion must be both nonlocal in prime
support and singular beyond every finite grading order.

\subsection{Finite vector transport}

The primitive-promoted CCM source maps each prime/depth cell to the actual
Gamma/co-Poisson interval coefficient before scalar orientation.  Exterior
multiplication yields a vector-valued chain map compatible with the
Calder\'on--Witt cup, incidence space, Koszul--Green section, and Green
retraction.

Cell labels cannot be fused away.  On the canonical shifted carrier, the
\((i,j)\)-boundary coordinate is exactly the reflection-odd Gamma
coefficient.  Distinct cell coordinates do not cancel.  For
\(\delta\ne0\), the canonical nonzero shifted carrier is outside the
admitted Hodge sector.  Consequently the existing finite packet
construction supplies no admission mechanism; a successful continuation
must remain relation-valued and preserve its cell labels.

\subsection{Hellinger residues and the off-diagonal boundary obstruction}

The reflected Hellinger/co-Poisson column has nonzero half-lattice
residues \(r_n\) at \(3/2+2n\).  For distinct primes \(p,q\), let
\(w_{p,q}\) be the normalized cross-prime incidence witness and define
\[
                         R_{p,q,n}=r_n w_{p,q}.
\tag{4.3}\label{eq:hellinger-incidence-residue}
\]
This is a nonzero actual Calder\'on--Witt incidence, and its normalized
incidence coordinate is exactly \(r_n\).  If \(\cA_a\) denotes the
Gamma--Hardy admitted sector at modular displacement \(a\), then
\[
 R_{p,q,n}\in\cA_a
 \quad\Longleftrightarrow\quad
 \mathfrak o_{\cA_a}(R_{p,q,n})=0
 \quad\Longleftrightarrow\quad a=0.
\tag{4.4}\label{eq:hellinger-relative-class}
\]
Moreover, its actual reflection-odd scalar boundary coordinate is
\[
 r_n\,\Gamma^{\mathrm{odd}}_{p,q,e}(a),
\tag{4.5}\label{eq:hellinger-odd-coordinate}
\]
and vanishes exactly at \(a=0\).

This closes the fixed-witness question negatively.  Nonzero residue
scaling cannot alter admission.  The obstruction is also polarized: the
existing vector Gamma transport has equal head and tail profiles before
Haar orientation, whereas the odd-Hodge image of the elementary head seed
has unequal channels at every positive source time.  Thus neither a fixed
incidence line nor a finite diagonal cell packet can furnish the missing
bridge.  A survivor must change the differential itself through an
off-diagonal, adjacent-grade coface whose infinite labelled chain has a
graph limit.

\subsection{Completed Xi-to-incidence bridge}

Fix distinct primes \(p,q\), and let \(w_{p,q}\ne0\) be the canonical
cross-prime Calder\'on--Witt incidence witness.  A finite Xi jet
\(a=(a_r)\) is realized by the actual incidence profile
\[
                         r\longmapsto a_r w_{p,q}.
\tag{4.6}\label{eq:finite-incidence}
\]
Pull back the Xi \(\ell^2\) norm along this coefficient-faithful map and
complete.  The resulting Hilbert space
\(\widehat{\cI}_{p,q}\) is canonically isometrically equivalent to the
completed Xi jet carrier.

There is a continuous Green comparison
\[
 U_{p,q}:H^1_{\!c}\longrightarrow\widehat{\cI}_{p,q}
\tag{4.7}\label{eq:completed-green}
\]
with the following unconditional properties:
\begin{align}
 \|U_{p,q}\ell\|&=\|\cG_\Xi\ell\|, \tag{4.8}\\
 \cl\ran U_{p,q}&=\widehat{\cI}_{p,q}, \tag{4.9}\\
 \operatorname{ev}_r(U_{p,q}\ell)
   &=\cG_\Xi(\ell)_r\,w_{p,q}. \tag{4.10}
\end{align}
On finite jets, \(U_{p,q}\) is the already constructed oriented
Calder\'on--Witt Green reflection.  Hence it is not a formal Hilbert copy
or a quotient that forgets coefficients.

Equations \eqref{eq:finite-incidence}--\eqref{eq:completed-green} close the
completed transport problem.  They do not define the correct closed
Gamma/Hodge admission at a divisor fibre.

\section{Relative admission-spectrum conservation}
\label{sec:admission}

\subsection{The route-independent law}

Let \(R\) be a field, let \(A_0,A_1\) be \(R\)-vector spaces, and fix
\[
 d:A_0\longrightarrow A_1,\qquad S\subset A_0,\qquad
 h:A_1\longrightarrow A_0.
\]
Here \(d\) is the source differential, \(S\) is the route's admitted
primitive sector, and \(h\) is any ambient primitive selector.  For a
family \(x:P\to A_1\), assume
\[
                           d\,h(x_p)=x_p.
\tag{5.A}\label{eq:primitive-selector}
\]

\begin{definition}[Relative admission spectrum]
The relative obstruction and its admission spectrum are
\[
 \mathfrak o_{S,d,h}(p)=[h(x_p)]
       \in Q_S(d):=A_0/(S+\ker d),
 \qquad
 \Sigma_S(d,h;x)=\{p:\mathfrak o_{S,d,h}(p)=0\}.
\tag{5.B}\label{eq:admission-spectrum}
\]
\end{definition}

Consider a second source complex \(e:B_0\to B_1\), an admitted sector
\(U\subset B_0\), and a chain square
\[
\begin{array}{ccc}
A_0&\xrightarrow{T_0}&B_0\\
{\scriptstyle d}\downarrow&&\downarrow{\scriptstyle e}\\
A_1&\xrightarrow{T_1}&B_1
\end{array}
\qquad
eT_0=T_1d,\qquad T_0(S)\subset U.
\tag{5.C}\label{eq:chain-square}
\]
\begin{theorem}[Relative admission-spectrum conservation]
\label{thm:main}
Under \eqref{eq:primitive-selector}, the following statements hold.
\begin{enumerate}[label=\textup{(\roman*)},leftmargin=2.2em]
\item \emph{Completeness.}
\[
 p\in\Sigma_S(d,h;x)
 \quad\Longleftrightarrow\quad
 \exists s\in S\ \text{such that}\ ds=x_p.
\tag{5.D}\label{eq:spectrum-completeness}
\]
\item \emph{Choice independence.}  If \(h'\) is another selector with
      \(dh'(x_p)=x_p\), then
      \(\mathfrak o_{S,d,h}(p)=\mathfrak o_{S,d,h'}(p)\).
\item \emph{Naturality.}  The square \eqref{eq:chain-square} induces
      \[
       \overline T_0:Q_S(d)\longrightarrow Q_U(e).
      \]
      If \(k\) selects primitives of \(T_1x_p\), then
      \[
       \overline T_0\bigl(\mathfrak o_{S,d,h}(p)\bigr)
        =\mathfrak o_{U,e,k}(p).
      \tag{5.E}\label{eq:obstruction-naturality}
      \]
\item \emph{Faithful conservation.}  If \(\overline T_0\) is injective,
      then
      \[
                    \Sigma_S(d,h;x)=\Sigma_U(e,k;T_1x).
      \tag{5.F}\label{eq:spectrum-conservation}
      \]
\item \emph{Exact loss law.}  If the target admits \(p\) and the source
      does not, then \(\overline T_0\) is noninjective and the specific
      nonzero class \(\mathfrak o_{S,d,h}(p)\) lies in
      \(\ker\overline T_0\).
\end{enumerate}
\end{theorem}

\begin{proof}
The class of \(h(x_p)\) vanishes precisely when
\(h(x_p)=s+n\) with \(s\in S\) and \(n\in\ker d\).  Applying \(d\)
gives \(ds=x_p\).  Conversely, if \(ds=x_p\), then
\(h(x_p)-s\in\ker d\), proving (i).  If \(h'\) is another selector, then
\(d(h(x_p)-h'(x_p))=0\); the two representatives therefore have the same
class, proving (ii).

For \(a=s+n\in S+\ker d\), sector preservation gives
\(T_0s\in U\), while the chain identity gives
\(eT_0n=T_1dn=0\).  Hence \(T_0a\in U+\ker e\), so \(T_0\) induces
\(\overline T_0\).  Both \(T_0h(x_p)\) and \(k(T_1x_p)\) are primitives
of \(T_1x_p\); their difference lies in \(\ker e\), proving
\eqref{eq:obstruction-naturality}.  Injectivity now makes source and target
vanishing equivalent, proving (iv).  If target vanishing occurs without
source vanishing, \eqref{eq:obstruction-naturality} places a nonzero source
class in \(\ker\overline T_0\), proving (v).
\end{proof}

The theorem does not identify unlike source spaces.  It states the exact
condition under which they may be identified: a bridge must commute with
their source differentials, preserve the admitted sector, and remain
injective after quotienting both admitted primitives and primitive
ambiguity.  Failure of admission conservation has only three possible
locations: failure of intertwining, loss of relative information, or a
singular change of domain.

\subsection{The Hedenmalm--Euler admission fracture}

Let \(H_a^x\) be the reflection-odd difference of the two Hedenmalm Borel
holonomy vectors at parameters \(ia\) and \(-ia\), based at \(x\).  Let
\[
 E_a(t)=\sqrt{1-e^{-t}}\,(e^{at}-e^{-at}),\qquad t>0,
\tag{5.G}\label{eq:weighted-euler-profile}
\]
be the polarized Euler--Hardy reflection profile with its actual tail
weight absorbed into the vector.

\begin{theorem}[Hedenmalm--Euler admission fracture]
\label{thm:hedenmalm-euler}
For every real basepoint \(x\), define
\[
 \Sigma_H(x)=\{a:H_a^x\in L^2(0,\infty)\},\qquad
 \Sigma_E=\{a:E_a\in L^2(0,\infty)\}.
\]
Then
\[
 \left(-\frac12,\frac12\right)\subset\Sigma_H(x),
 \qquad \Sigma_E=\{0\}.
\tag{5.H}\label{eq:analytic-spectra}
\]
In particular, for \(0<|a|<1/2\), no linear raw-profile map \(T\) can
satisfy both
\[
                  T H_a^x=E_a,\qquad
                  H_a^x\in L^2\Longrightarrow T H_a^x\in L^2.
\tag{5.I}\label{eq:no-regular-intertwiner}
\]
No boundedness, injectivity, inverse, or reverse domain implication is
needed.
\end{theorem}

\begin{proof}
Both Hedenmalm holonomy vectors are square integrable when \(|a|<1/2\),
and \(L^2\) is closed under subtraction; hence the first inclusion.  The
squared norm of \eqref{eq:weighted-euler-profile} is
\[
 \int_0^\infty(1-e^{-t})|e^{at}-e^{-at}|^2\,dt,
\]
which is finite exactly when \(a=0\), giving the second identity.  If a
map in \eqref{eq:no-regular-intertwiner} existed at nonzero \(a\), the
first spectrum would put \(H_a^x\) in \(L^2\); domain preservation and
intertwining would then put \(E_a\) in \(L^2\), contradicting
\(\Sigma_E=\{0\}\).
\end{proof}

This is the first direct cross-route consequence.  A regular common
Hilbert carrier cannot merge the permissive Borel source with the sharp
Euler--Hardy boundary.  The missing bridge is necessarily a singular
domain restriction before scalarization.  The result does not show that
Xi zero fibres land in that singular domain.

\subsection{Absolute orientation as one incidence specialization}

Let \(A^1\) be the all-prime finite-support Calder\'on--Witt degree-one
source, \(A^2\) its first suspension, and
\[
 \smile:\Lambda^2A^1\longrightarrow A^2,\qquad
 I=\ker(\smile).
\]
The incidence inclusion \(\iota:I\hookrightarrow\Lambda^2A^1\), exterior
cup, Green retraction \(h\), and cup section form a contraction.  Thus the
unoriented complex is acyclic.  Its normalized Haar orientation trace
\(\tau:A^2\to\C\) defines
\[
 D x=\bigl(\smile x,-\tau(\smile x)\bigr),\qquad
 T(z,c)=\tau(z)+c .
\tag{5.1}\label{eq:oriented-complex}
\]
One has \(\operatorname{im}D=\ker T\), and the resulting top cohomology is
one-dimensional.

\begin{theorem}[All-prime orientation--transgression completeness]
\label{thm:orientation}
The all-prime oriented incidence source has the following three
properties.
\begin{enumerate}[label=\textup{(\roman*)},leftmargin=2.2em]
\item Every linear consumer \(F:A^2\oplus\C\to V\) has a canonical
      decomposition
      \[
                         F=F_{\mathrm{abs}}+F_{\partial},
      \tag{5.2}\label{eq:observer-splitting}
      \]
      where \(F_{\mathrm{abs}}\) is boundary-blind and
      \(F_{\partial}\) vanishes on the orientation line.  This is the
      unique decomposition with those two properties.
\item The following are equivalent:
      \[
      F_{\partial}=0
      \quad\Longleftrightarrow\quad F\circ D=0
      \quad\Longleftrightarrow\quad
      \exists!\,\phi:\C\to V\ \text{such that}\ F=\phi\circ T.
      \tag{5.3}\label{eq:observer-factorization}
      \]
\item Let \(S\subset I\) be any admitted incidence sector and let
      \(x\in\Lambda^2A^1\) satisfy \(\smile x=0\).  Its relative
      transgression class
      \[
                         \mathfrak o_S(x)=[h x]
                         \in I/(S+\ker\iota)
      \]
      obeys
      \[
                         \mathfrak o_S(x)=0
                         \quad\Longleftrightarrow\quad
                         h x\in S.
      \tag{5.4}\label{eq:relative-transgression}
      \]
\end{enumerate}
\end{theorem}

\begin{proof}
The Green decomposition of every top cochain is
\[
                       z=D(H z)+\jmath(Tz).
\]
Applying \(F\) defines
\(F_{\partial}(z)=F(D(H z))\) and
\(F_{\mathrm{abs}}(z)=F(\jmath(Tz))\), proving the split.  Its vanishing
laws and the Green decomposition give uniqueness.  If \(F\circ D=0\),
then \(F_{\partial}=0\) and \(F(z)=F(\jmath(Tz))\), so \(F\) factors
through \(T\).  The identity \(T\circ\jmath=1\) gives uniqueness; the
converse follows from \(T\circ D=0\).  Finally, the contraction identifies
the complete relative obstruction with the class of \(hx\).  Since
\(\iota\) is injective, \(\ker\iota=0\), so this class vanishes exactly
when \(hx\in S\).
\end{proof}

Equations~\eqref{eq:observer-splitting}--\eqref{eq:observer-factorization}
are the absolute/relative observation law: every regular consumer has one
orientation compression and one exact boundary-sensitive remainder.
Equation
\eqref{eq:relative-transgression} is the relative law: every further
invariant is sector-sensitive primitive escape.  It does not identify
the admitted sectors of different routes.  It states exactly what each
route must preserve before passing to its consumer.

For any parameter family \(w_p\) of all-prime exterior cocycles, the
incidence instance of Theorem~\ref{thm:main} is therefore the exact set
identity
\[
 \Sigma_S(w)=
 \left\{p:h(w_p)\in S\right\}
 =\left\{p:\exists s\in S\text{ with }\iota(s)=w_p\right\}.
\tag{5.J}\label{eq:deninger-spectrum}
\]
This is not a metaphorical comparison with the analytic spectra in
\eqref{eq:analytic-spectra}; all three are vanishing spectra of the same
relative-primitive construction.  Their different shapes record different
admitted sectors, not different choices of Green operator.

\subsection{The completion threshold}

\begin{lemma}[Dense closed-admission threshold]
\label{thm:dense-admission}
Let \(U:E\to H\) have dense range in a topological space \(H\), let
\(d\in H\), and let \(P\) be any oriented target condition.  Define a
closed admission realization to be a closed set \(K\subset H\) satisfying
\[
                U(E)\subset K,\qquad d\in K\Longleftrightarrow P.
\tag{5.5}\label{eq:abstract-admission}
\]
Then a closed admission realization exists if and only if \(P\) holds.
\end{lemma}

\begin{proof}
If \(K\) exists, closedness and density give
\(K\supset\cl U(E)=H\); hence \(d\in K\) and \(P\) follows.  Conversely,
if \(P\) holds, \(K=H\) satisfies \eqref{eq:abstract-admission}.
\end{proof}

This elementary closure lemma is not the paper's universal law.  Its sole
role is to close the completed-incidence specialization once a route has
constructed its dense source transport, detector, and exact trace without
assuming \(P\).

\subsection{Completed-incidence specialization}

Fix a divisor coordinate \(r\) and define
\[
 e_{r;p,q}:=U^{\mathrm{jet}}_{p,q}(\mathbf1_r)
             \in\widehat{\cI}_{p,q},
\tag{5.6}\label{eq:canonical-vector}
\]
where \(\mathbf1_r\) is the unit finite jet.  Decoding at \(r\) gives
\[
                  \operatorname{ev}_r(e_{r;p,q})=w_{p,q}.
\tag{5.7}\label{eq:canonical-decodes}
\]

Let \(\cA_\delta\) be the algebraic Gamma/co-Poisson Hodge-admitted
sector, the kernel of the actual reflection-odd boundary defect.  The
canonical witness detects displacement:
\[
                       w_{p,q}\in\cA_\delta
                       \quad\Longleftrightarrow\quad
                       \delta=0.
\tag{5.8}\label{eq:witness-detects}
\]

\begin{definition}[Completed-incidence admission]
\label{def:admission}
A completed-incidence admission at displacement \(\delta\) is a closed
set \(K_\delta\subset\widehat{\cI}_{p,q}\) such that
\[
 U_{p,q}(H^1_{\!c})\subset K_\delta
\tag{5.9}\label{eq:source-landing}
\]
and whose one-coordinate finite trace is exact:
\[
 e_{r;p,q}\in K_\delta
 \quad\Longleftrightarrow\quad
 w_{p,q}\in\cA_\delta.
\tag{5.10}\label{eq:finite-trace}
\]
\end{definition}

No linearity, bounded boundary operator, or equality on the whole finite
core is required.

\begin{corollary}[Completed-incidence source-admission law]
\label{thm:xi-admission}
Fix distinct primes \(p,q\) and one divisor coordinate \(r\).
\begin{enumerate}[label=\textup{(\roman*)},leftmargin=2.2em]
\item For every real displacement \(\delta\), a completed-incidence
      admission exists if and only if \(\delta=0\).
\item At a nontrivial divisor point \(\rho\), admission at
      \(\delta_\rho=2(\Rea\rho-\tfrac12)\) exists if and only if
      \(\Rea\rho=\tfrac12\).
\item Admissions exist at every nontrivial divisor point if and only if
      the Riemann hypothesis holds.
\end{enumerate}
\end{corollary}

\begin{proof}
Apply Theorem~\ref{thm:dense-admission} to the dense map \(U_{p,q}\), the
detector \(e_{r;p,q}\), and the target condition
\(w_{p,q}\in\cA_\delta\).  Equation~\eqref{eq:witness-detects} gives
(i).  Substitution of \eqref{eq:delta} gives (ii), and quantifying (ii)
over the nontrivial divisor gives (iii).
\end{proof}

\begin{remark}[Why closedness is a threshold, not a universal starting point]
If admission is the kernel of a closed possibly unbounded boundary
operator, it is closed and Corollary~\ref{thm:xi-admission} applies.  A
dense nonclosed Hedenmalm radical can instead exclude the decisive
completed detector.  Product-topology closure can become maximally
vertical and output-free.  Other routes have not yet constructed a dense
transport into their intended admission carrier.  Closedness therefore
separates one completed route from these different earlier obstructions;
it does not erase them.
\end{remark}

\section{Route atlas: distinct open cores under one law}
\label{sec:routes}

Theorem~\ref{thm:main} does not collapse all constructions into the
completed-incidence model.  It says that a route must preserve a
source-relative Green primitive before its scalar consumer.  The
following atlas records where each route currently stops.

\begingroup
\footnotesize
\setlength{\tabcolsep}{4pt}
\renewcommand{\arraystretch}{1.18}
\begin{longtable}{
  >{\raggedright\arraybackslash}p{0.16\linewidth}
  >{\raggedright\arraybackslash}p{0.28\linewidth}
  >{\raggedright\arraybackslash}p{0.50\linewidth}}
\toprule
\textbf{Route} & \textbf{Unconditional prefix} &
\textbf{First unproved producer theorem}\\
\midrule
\endfirsthead
\toprule
\textbf{Route} & \textbf{Unconditional prefix} &
\textbf{First unproved producer theorem}\\
\midrule
\endhead
\midrule
\multicolumn{3}{r}{\emph{Continued on next page}}\\
\endfoot
\bottomrule
\endlastfoot
C1: singular Hedenmalm
& Genuine Borel \(L^2\) source throughout the centered half-strip;
  full-source positive-form classification; exact admission-spectrum
  mismatch with the Euler--Hardy carrier
& Construct a source-selected domain-singular boundary form retaining
  divided-zero levels and intertwining the sharp polarized profile\\

E: M\"obius
& Exact zero/one all-prime boundary factorization; renewal kernel; finite
  conductance and Hall-cut laws; local hard-support producers excluded
& Prove critical strict-half cancellation across smooth/one-large-prime
  seam\\

G: Suzuki
& Finite all-angle Hermite--Biehler carrier; exact first-defect cone,
  endpoint budgets, and safe-shift Hankel gate
& Construct pole-compatible two-angle meromorphic continuation with
  source-selected nondegeneration and Xi/de Branges identification\\

J: Julia relative current
& Compound-Poisson source, exact finite capacity hierarchy, cofinal
  depth-nonlocal reflected countertransport, bounded pretrace capacity
& Construct source-relative Hodge polarization controlling negative
  first-clock excursion before singular codiagonal, or prove a
  class-complete no-go\\

L: Calder\'on--Witt/Hodge
& All-prime oriented source complex; finite Gamma/co-Poisson boundary
  comparison; dense norm-preserving completed Xi incidence bridge;
  fixed-witness Hellinger residue classified exactly
& Construct an adjacent-grade off-diagonal Hellinger/Beta coface whose
  labelled cofinal chain is graph-Cauchy and lands the Xi Green primitive\\

N: prime-normal domain
& Exact decomposition of prime-normal energy into vertical clock and
  reflected radial energy
& Prove uniform sub-\((\log X)^2\) estimate on the closed theta Grushin
  source domain with full Jordan tower\\

O: Connes/CCM
& Literal restricted-Weil matrices; exact symmetry and clock
  displacement; prolate Rayleigh ground-lock classification
& Prove cofinal restricted-Weil floor convergence from literal
  prime/Gamma/prolate source\\

P: Thorin/Markov
& Compact moment determinacy and uniqueness of every possible positive
  divisor carrier
& Construct the compact positive carrier from source without divisor
  input\\

Q: theta--Cassini
& Positive theta likelihood carrier, all raw Hankel blocks, nonlinear
  response, unique calibrated reciprocal exponent
& Control calibrated carrier from degree two onward or prove boundary
  Jensen/jet saturation\\

R: primitive current
& Wold activation, valuation disintegration, target-free physical
  M\"obius--theta source, exact complex source-to-Xi landing,
  Hellinger/Gamma coordinates, primitive promotion, Julia/KMS resolution
& Complete the singular prime-normal source graph and prove its
  source-derived Hellinger/Ward admission before scalarization\\

analytic--Hurwitz
& Completed-density relation at each divisor point and corresponding
  Hurwitz circuit exactly on the critical line
& Construct Euler--Gamma--theta relation comparison into the circuit\\

HT/recurrence
& Hypertranscendental facet and residue sources; balanced projective
  relation lattice with exact rank laws
& Construct arithmetic divisor reporter and prove relation protection
  under the source-allowed gauge or perturbation class\\
\end{longtable}
\endgroup

\subsection{Why these remain different routes}

Rows E, N, O, P, and Q lack a source-domain estimate or source-built
carrier.  Route C1 has a genuine Hilbert source but, by
Theorem~\ref{thm:hedenmalm-euler}, any comparison with the sharp
Euler--Hardy carrier must be domain-singular.  Rows G,
analytic--Hurwitz, and HT/recurrence lack a faithful continuation into
their target relation.  Route R now has physical source-to-Xi landing but
lacks the singular prime-normal-to-Hodge admission.  Route J has constructed a
cofinal vector transport but lacks the source-relative polarization which
survives its singular scalar recovery.  Route L alone has a dense
completed incidence carrier and an exact finite detector, so
Corollary~\ref{thm:xi-admission} sharpens its remaining theorem to a
source-derived off-diagonal cofinal graph landing.  These are not
interchangeable gaps.

Terminalized first-clock and Gamma--Jacobi reader coordinates, and
determinant, Calkin, Sonin, Ward, or recurrence equivalences without an
upstream producer, are not additional routes.  They are consumers or
closed interfaces.  Conversely, retirement of a bounded or finite
submechanism does not retire an open route whose singular source category
survives.

\subsection{The analytic--monodromy version of admission}

At every divisor point, the theta completed-density source has a relation
\[
                         E_\rho+E_\rho^\#=0.
\tag{6.1}\label{eq:analytic-relation}
\]
The corresponding Hurwitz pair is a differential circuit exactly when
\(\Rea\rho=\tfrac12\).  Therefore
\[
 \text{every relation \eqref{eq:analytic-relation} is preserved}
 \quad\Longleftrightarrow\quad \mathrm{RH}.
\tag{6.2}\label{eq:circuit-rh}
\]
The missing theorem in that route is not \eqref{eq:circuit-rh}; it is the
producer comparison which transports the represented source relation into
the Hurwitz circuit.  This is the algebraic analogue of
\eqref{eq:source-landing}.

\subsection{The Frobenius--Hodge version of admission}

The normalized arithmetic Real structure commutes with Frobenius
unconditionally.  Positive Hodge duality factors as that source Real
structure followed by the arithmetic Weil transport
\(\rho\mapsto1-\overline\rho\).  Frobenius commutes with the Weil
transport exactly under RH.  Thus the former phrase ``weight-one
conservation'' has a precise meaning:
\[
              [\Theta,\cW_{\mathrm{Weil}}]=0.
\tag{6.3}\label{eq:frobenius-hodge}
\]
The completed admission theorem asks for a source construction of this
intertwining, not another proof that its vanishing is equivalent to RH.

\section{No-go and positive theorems as shadows}
\label{sec:shadows}

The source-admission law explains why previously separate results recur.

\subsection{Dense radicals and maximal verticality}

A dense nonclosed Hedenmalm radical can contain every finite source test
without containing a decisive completed vector.  Closing it forces loss of
orientation information.  The prime-normal product completion is the
relation-valued extreme: its vertical part becomes the whole target.
In admission-spectrum language, the first changes the domain on which the
quotient class is defined; the second kills every quotient class.  One is
necessarily singular, while the other is maximally noninjective.

\subsection{Wold--Poisson ancestry and Rees superselection}

Wold shell depth and prime-normal codimension are independent filtrations
of one integer source.  At nonzero modular displacement their reflected
valuation sectors acquire different asymptotic species.  A finite-Rees
operator preserves the flat radical, so it cannot produce an admission
set with the finite trace \eqref{eq:finite-trace}.  The singular growth in
\eqref{eq:beyond-rees} is the cost of leaving one admission-spectrum
class for another without killing the source quotient.

\subsection{Sonin, Calkin, and positive Ward consumers}

Sonin determinants and Calkin compactness detect whether a transported
source remains in a regular operator ideal.  Four-unit Ward positivity
measures the squared transverse defect.  These are stable consumers of
closed admission, but they do not preserve the unsquared finite trace by
themselves.  Off-line quartets can carry positive \(8\pi\)-quantized Ward
mass; positivity counts the defect rather than annihilating it.

\subsection{Hypertranscendence and recurrence}

Differential independence prevents any finite jet or finite
differential-algebraic sector from manufacturing the full source
relation.  Projective recurrence gives the complementary rank statement:
observable temporal rank plus exact relation rank is conserved.  Under
unrestricted perturbations no nonzero relation is protected.  Both facts
say that admission must be source-relative.  A relation is meaningful
only together with the perturbations or gauges that the source allows.

\subsection{The common law}

All these shadows can be summarized without identifying their categories:
\[
\boxed{
\begin{minipage}{0.84\linewidth}
\centering
Relative primitive classes are independent of Green choice and natural
under source chain maps.  Every injective admitted-quotient transport
conserves their vanishing spectrum.  A spectral mismatch forces a precise
kernel class or a singular change of source domain.
\end{minipage}}
\tag{7.1}\label{eq:universal-law}
\]
Theorem~\ref{thm:main} proves this statement.  The
orientation--transgression split is its all-prime incidence shadow;
Theorem~\ref{thm:hedenmalm-euler} is its analytic cross-route
consequence.  Corollary~\ref{thm:xi-admission} proves that closed admission
is exactly the critical-line condition for the completed-incidence
instance.  None of these results supplies source-derived global landing.

\section{What remains for a full RH proof}
\label{sec:proof}

\subsection{Common closed components}

The following parts no longer require conjectural objects, although no
single route uses all of them.
\begin{enumerate}[leftmargin=2.2em]
\item The signed divisor displacement and full multiplicity-weighted
      outer charge have source realizations.
\item Mass-one prime-normal, valuation, CCM crease, Xi jet, theta
      relation, hypertranscendental, and recurrence sources exist.
\item The Hellinger/KMS polarization, Beta/Gamma transform, transverse
      four-unit source vector, finite Gamma/co-Poisson transport,
      Calder\'on--Witt incidence, and Green retraction are explicit.
\item Target-free physical M\"obius--theta rungs have exact
      source-to-Xi Fourier--Laplace landing on their honest convergence
      domain.
\item The completed Xi Green source maps coefficient-faithfully and
      densely into an honest incidence completion.
\item One finite completed incidence coordinate decodes to the actual
      nonzero Gamma witness.
\item That witness is Hodge-admitted exactly at zero displacement.
\item Every Hellinger--Beta half-lattice residue has a nonzero actual
      incidence lift whose relative Gamma class vanishes exactly at zero
      displacement.
\item Hedenmalm and Euler--Hardy reflection sources have incompatible
      finite-energy admission spectra; every intertwining bridge between
      them is necessarily domain-singular.
\item Regular product completion, finite-Rees sewing, bounded
      label-blind polarization, local positive carriers, bounded scalar
      codiagonals, and finite differential closure cannot replace a
      source-relative oriented transgression.
\end{enumerate}

\subsection{The completed-incidence route}

Fix any distinct primes \(p,q\) and any one divisor coordinate \(r\).
For Route L, the terminal criterion is the source-derived construction of
\[
\boxed{
\begin{minipage}{0.84\linewidth}
For every \(\rho\in\Div(\Xi)\), construct a closed set
\(K_\rho\subset\widehat{\cI}_{p,q}\) containing
\(U_{p,q}(H^1_{\!c})\), and prove the single finite trace identity
\begin{center}
\(\displaystyle
 e_{r;p,q}\in K_\rho
 \quad\Longleftrightarrow\quad
 w_{p,q}\in\cA_{\delta_\rho}.
\)
\end{center}
\end{minipage}}
\tag{8.1}\label{eq:missing}
\]
Corollary~\ref{thm:xi-admission} then gives RH immediately.

The set in \eqref{eq:missing} cannot be chosen after inspecting
\(\delta_\rho\); that would merely restate the equivalence.  A valid
producer must define \(K_\rho\) uniformly as the graph domain or kernel of
the source-pointed Gamma/Hellinger boundary and derive Green landing from
the prime-power source before divisor scalarization.  In the existing
boundary coordinate, this requires the labelled finite-tail estimate
\[
 t\longmapsto |A(i,j)|^2(1-e^{-t})
 \left|e^{\delta_\rho t}-e^{-\delta_\rho t}\right|^2
 \in L^1(0,\infty)
\tag{8.2}\label{eq:tail-energy}
\]
for every nondegenerate source cell, proved from source data rather than
assumed at the divisor.  Existing results classify
\eqref{eq:tail-energy} as equivalent to zero displacement; they do not
prove it.

The residue computation rules out treating \(K_\rho\) as an arbitrary
closure or as the completion of a fixed witness line.  Each nonzero
Hellinger residue remains outside the admitted sector at every nonzero
displacement, and every current finite Gamma packet is diagonal in the
head/tail polarization.  The surviving producer is consequently more
specific:
\[
\boxed{
\begin{minipage}{0.84\linewidth}
Construct an adjacent-grade Hellinger--Beta coface on the labelled
prime-power incidence chain.  Its finite boundaries must telescope, its
partial Green primitives must be Cauchy in the closed boundary graph norm,
and its graph limit must be the completed Xi Green primitive.  Coefficients
and domain must be derived from the physical source before a divisor point
is selected.
\end{minipage}}
\tag{8.3}\label{eq:missing-cofinal-coface}
\]
This formulation identifies the required transport and excludes
fixed-witness, finite-cell diagonal, bounded-Hilbert, and arbitrary
closed-set repairs.  The construction in
\eqref{eq:missing-cofinal-coface} is not presently available.

The topological formulation is nevertheless strictly weaker than earlier
completion targets.  It does not ask for:
\begin{enumerate}[leftmargin=2.2em]
\item a bounded Gamma/Hodge defect on the completion;
\item equality of completed and algebraic kernels on every finite vector;
\item annihilation proved separately on the whole incidence space;
\item a new Hilbert space, source, positivity form, or divisor
      equivalence.
\end{enumerate}
Closedness, dense source landing, and one correct oriented trace suffice.

\subsection{Independent sufficient theorems in the other routes}

The other routes do not first need to construct
\(K_\rho\subset\widehat{\cI}_{p,q}\).  Each has its own sufficient theorem
and an already constructed connector to zero exclusion.
\begin{enumerate}[leftmargin=2.2em]
\item \emph{Singular Hedenmalm:} construct the source-selected singular
      boundary domain which maps the permissive Borel spectrum to the
      sharp Euler--Hardy spectrum while retaining divided-zero levels.
\item \emph{Critical M\"obius:} prove the strict-half zero/one cancellation
      across the macroscopic smooth/one-large-prime seam.
\item \emph{Suzuki:} construct the pole-compatible two-angle meromorphic
      producer, prove source-selected nondegeneration, and identify its
      limit with the Xi/de Branges target.
\item \emph{Principal-relative Julia:} prove a source-Hodge polarization
      which bounds negative first-clock excursion by the bounded pretrace
      capacity before singular codiagonal recovery.
\item \emph{Prime-normal domain:} prove the uniform source-domain
      sub-log-square estimate on the closed theta Grushin domain.
\item \emph{Connes/CCM:} prove that the true cofinal restricted-Weil
      spectral floor tends to zero from the literal prime/Gamma/prolate
      source.
\item \emph{Thorin/Markov:} construct the compact positive carrier whose
      uniqueness is already known.
\item \emph{Theta--Cassini:} control the calibrated positive carrier from
      degree two onward, or prove the boundary Jensen/jet saturation
      inequality directly.
\item \emph{Primitive current:} extend the constructed physical
      M\"obius--theta source-to-Xi landing to its singular prime-normal
      graph and prove source-derived Hellinger/Ward admission before
      scalarization.
\item \emph{Analytic--Hurwitz:} transport the theta completed-density
      relation into the Hurwitz differential circuit.
\item \emph{Hypertranscendental/recurrence:} identify an arithmetic
      divisor reporter and prove that its relation survives precisely the
      gauges or perturbations admitted by the source.
\end{enumerate}

Any one of these theorems, with its proved connector, would settle RH.
None is presently proved.  Theorem~\ref{thm:main} explains their shared
architecture: each must preserve a relative primitive class which its
downstream consumer cannot reconstruct.  It does not turn them into one
technical problem.

\section{Scope and conclusion}

This paper proves a route-independent conservation law and a concrete
cross-route singularity theorem, not RH.  It neither assumes nor
constructs the missing source estimates, singular continuations, or
admitted graph domains in Section~\ref{sec:proof}.  In the
completed-incidence instance, closed realization exists exactly on the
critical line and global realization is equivalent to RH.

The universal law is:
\[
\boxed{
\begin{minipage}{0.84\linewidth}
\centering
The vanishing spectrum of a relative source primitive is independent of
primitive choice and is conserved by every injective admitted-quotient
chain transport.  If two route spectra differ, every intertwining bridge
must kill a nonzero relative class or change source domain singularly.
\end{minipage}}
\tag{9.1}\label{eq:conclusion-law}
\]
\end{document}
